\begin{abstract}
Model merging has emerged as a dominant, training-free paradigm for fusing multiple task-specific neural network experts into a single unified multi-task model. However, direct linear addition of fine-tuned weights, as in standard \emph{Task Arithmetic}, typically suffers from catastrophic representational collisions and severe interference in weight space, causing a significant collapse in performance. While existing approaches attempt to resolve these conflicts using complex sign consensus heuristics or stochastic dropout masks, we propose \textbf{Sparsity-Guided Task Arithmetic (SG-TA)}, a remarkably simple and deterministic weight-space regularization framework. SG-TA decouples the sparsification phase by applying magnitude-based binary masking to individual task-specific update vectors before merging, filtering out harmful orthogonal parameter noise and destructive interference. Furthermore, we explore two masking scopes: \emph{Global Quantile (GQ)} masking and \emph{Layer-wise Quantile (LQ)} masking. To optimize the merging hyperparameters without overfitting, we leverage \emph{Offline Few-Shot Validation Tuning (OFS-Tune)} using only 10 samples per task. Adhering to a strictly empirical research philosophy, we conduct an exhaustive, statistically rigorous evaluation across multiple random calibration seeds on a 4-dataset multi-domain benchmark (MNIST, FashionMNIST, CIFAR-10, SVHN) using a Vision Transformer (ViT-Tiny) backbone. Our results show that SG-TA under global quantile masking (GQ) achieves an average joint mean accuracy of \textbf{61.40\% $\pm$ 1.39\%} (a $+15.08\%$ absolute improvement over Naive Uniform Task Arithmetic and $+2.17\%$ over Optimized Task Arithmetic), outperforming more complex baselines like TIES-Merging (60.64\% $\pm$ 1.30\%) and DARE-Merging (58.44\% $\pm$ 3.02\%), and layer-wise scaling without pruning (32.44\% $\pm$ 5.49\%). This provides strong empirical validation that spatial parameter regularization via simple, global magnitude-based weight masking is highly effective for model merging.
\end{abstract}
